\chapter{公平性}
\label{chap:trust-fairness}

一个培训名额分配模型在总体上预测准确，却可能把大量合格申请者集中拒绝在某个群体中。判断这种结果是否公平，既需要知道错误如何分布，也需要知道名额应按什么原则分配。即使所有群体的准确率相同，误拒与误收的比例、获得机会的数量和实际后果仍可能不同。

前两章讨论预测的可靠性与变化条件下的鲁棒性。这些分析能够揭示错误集中出现的条件，却还没有规定不同人之间怎样分配机会和负担。名额分配的公平性需要明确自己的判断原则，不能由预测准确或性能稳定直接推出。本章要回答的是：分配系统应保护什么、怎样判断保护要求是否满足，以及怎样改变造成不合理差异的环节。概念部分先说明公平的含义、不公平的表现及形成原因；判定部分把保护维度、准则和检验方法对应起来；干预部分则把公平要求落实到数据、训练和决策流程。判定准则决定干预目标，形成机制决定措施选择；措施是否有价值，还要看实际效果和长期反馈，不能只看某个指标是否下降。

\section{模型的公平性}
\label{sec:fairness-concepts}

\subsection{公平性的含义}
\label{sec:fairness-meaning}

\term{公平性}（fairness）要求模型及其使用过程依据任务认可的原则，对不同个体和群体合理分配机会、收益、错误与负担。它既关心谁得到什么，也关心依据什么作出决定，以及受影响者能否获得恰当处理。公平不等于所有人的输出相同；关键是哪些差异有合理依据，哪些差异违背了事先说明的保护要求。

例如，培训名额可以按现有资格、培训后的预期收益或弥补机会不足的需要分配。这些原则可能产生不同名单，不能靠准确率选出唯一正确原则。公平要求包含规范性判断，即在当前任务中什么差异应被允许、什么差异应被纠正。统计方法可以检查已经明确的要求，却不能单独决定应采纳哪一种分配原则。

设$\bm{x}$为可用特征，$a$为需要关注的群体属性，$y\in\{0,1\}$为任务标签，$\hat y$为二元决策。本章在名额分配例子中令$y=1$表示按既定标准合格，$\hat y=1$表示获得机会。标签的这一含义必须在实际任务中核查：历史上是否被录取，不等于是否具备资格。邮件例子中的正决策是隔离，本章则是提供机会；同一个错误率在不同任务中可能对应不同后果，不能脱离标签含义判断越高或越低越好。

判断公平还需要限定系统边界。数据中不同群体的数量、模型输出的分数、实际录取决定、参加培训的成本和最终收益，是不同层面的对象。预测概率相同，可能因决策规则不同而产生不同机会；录取决定相同，也可能因费用或服务条件不同而产生不同实际收益。本章把模型放回使用过程讨论，避免用某一个环节的均等代替对完整后果的判断。

\subsection{不公平的表现}
\label{sec:fairness-manifestations}

可以从具体的待遇与后果识别需要调查的问题。设两个群体各有20名按同一标准认定的合格申请者，其中一组18人获选，另一组只有10人获选。在这个教学例子中，若保护要求是合格者不因群体不同而失去机会，那么两组不同的误拒程度就是需要解释和纠正的现象。若只知道两组总体录取率不同，却不知道采用了什么资格标准，就还不能直接断定哪一方受到不合理处理。

另一种问题发生在个体之间。假设两名申请者在所有被认可的任务条件上相同，系统却只因与任务无关的地区标签而稳定地给予不同录取机会，那么即使两个地区的总体录取率相同，也没有消除这种差别待遇。相反，对不同资格者作出不同决定是否合理，需要回到资格标准本身，而不能仅按输出是否一致判断。

图~\ref{fig:trust-fairness-opportunity}用报名材料的装饰差异代替地区标签，突出比较所需的条件。假定材料中的三个资格项目已经覆盖当前认可的全部分配依据，且没有名额抽签等允许的随机机制，那么仅因封面不同而稳定地产生不同待遇，就违反了这一情景的个体一致性要求。现实审计需要验证这些假定，不能只凭两份材料外观相近便得出结论。

\begin{figure}[!htb]
  \centering
  \bookdraftgraphic{figures/trust-fairness-opportunity.png}
  \caption{相同资格下的不同待遇。上下两行表示同一系统对两名申请者的处理；三个相同图标表示相同的认可资格，条纹与圆点仅是无关的封面装饰。一人获得培训机会，另一人被拒，提示需要调查的差别待遇。图示为构造情景。}
  \label{fig:trust-fairness-opportunity}
\end{figure}

还需区分\term{分配性伤害}（allocative harm）与\term{表征性伤害}（representational harm）。前者体现在谁得到名额、服务或资金，后者体现在谁被刻板描述、忽视或贬低。培训推荐系统可能给予各组相同数量的推荐，却在没有事实依据时反复把某些人描述为不适合技术岗位；机会计数相同，并未覆盖这种表达带来的影响。另一个需要检查的现象是服务负担：某组即使最终获选比例相同，却需要更多证明、更长等待或更多申诉才能完成申请。

实际研究也展示了总体指标掩盖差异的情形。Buolamwini和Gebru的Gender Shades按研究中的肤色与性别标注交叉分析商用图像性别分类器，发现不同交叉组的错误水平存在明显差异。这项工作说明为何需要拆分评价；其结论限定在研究样本、标注方式和当时系统上，不能把肤色当作种族的无误替代，也不能把图像性别标注等同于人的完整性别身份\scite{buolamwini2018gendershades}

上述表现分别指向机会与负担、预测质量、个体待遇和社会表征。它们帮助确定要检查什么，还不能单独回答差异由哪个环节造成；同一种表现可能来自不同原因，因而需要不同干预。

\subsection{不公平的形成原因}
\label{sec:fairness-causes}

问题可能从任务目标开始。如果把历史录取结果直接当作“是否合格”的标签，模型就可能学习过去的筛选习惯，而非当前认可的资格。即使模型准确复现了标签，也未必符合新的公平要求。这类问题需要重新说明预测目标和决策用途，不能只增加训练样本。

数据采集与测量会改变模型看见的世界。只从线上报名者收集样本，可能遗漏不易使用线上渠道的人；不同语言版本的测验难度不同，可能把测量差异写成能力差异；标注者使用不一致标准，又可能造成分组标签误差。样本不足、测量不等价和标签偏差是三种不同问题，补数量只能直接缓解其中一部分。

学习过程也可能放大已有差异。若训练目标只最小化总体平均错误，小群体上的较大损失可能被多数样本的改善抵消。某些特征还会传递被删除属性的信息，例如地区与既有资源高度相关时，删除显式群体字段并未消除这种关联。是否允许利用某项信息，需要依据任务原则及其影响路径判断，不能把“看不到群体字段”当成公平证明。用于审计、用于训练约束与直接用于决策的群体信息，也应分别说明用途。

部署环节决定分数怎样变成实际待遇。统一阈值未必对应各组相同错误率，不同阈值也未必符合个体一致性要求；人工复核、收费方式与等待时间还会改变模型以外的负担。只调整预测器而保留有问题的服务流程，可能使差异继续存在。

决策还会改变未来数据。在名额分配中，只有获得培训者才有后续成绩，形成\term{选择性标签}（selective labels）：结果能否观测取决于既有决策。如果直接用这些成绩训练，模型可能持续缺少被拒者潜在表现的证据。机会分配还可能改变下一轮的能力和资源，从而形成反馈。纠正这些问题需要额外观测或可识别假设，不能仅在已有记录上更换一个公平指标。

形成原因用于定位干预环节，公平判定则用于判断哪些要求没有满足。例如，在其他输入不变时修改邮政编码，若预测随之改变，可以得到模型依赖该输入的线索，却还不能判断这种依赖是否合理；预测没有改变，也不排除其他代理变量传递相同信息。后文第~\ref{sec:interpretability-local}节将介绍更系统的行为解释方法。需要把机制线索与明确的保护要求、比较对象和证据结合起来。

\section{公平性的判定}
\label{sec:fairness-criteria}

判定公平性要把三个层次连起来：判定维度说明关心哪一种待遇或后果，公平准则规定什么关系应当成立，检验方法则用数据或实验检查该关系。例如，关心合格者是否同样有机会，是判定的内容；要求各组真正例率相同，是具体准则；按群体统计合格者中的获选比例并估计差距，是检验方法。三者不能只用一个算法名称代替。

下面从机会与负担的分配、预测错误与分数含义、相似个体的处理、决策依据与影响路径、生成内容中的社会表征五个方面展开。前两个方面主要通过群体统计检验，第三个通过个体成对比较，第四个需要关于影响路径的证据，第五个需要内容评价与受控比较。每个方面都说明比较对象、准则、操作方法与局限。末尾再讨论如何综合这些证据形成结论；准则取舍和证据强度是所有方面共同面对的问题，不是另外两种公平维度。

\subsection{机会与负担的分配}
\label{sec:fairness-group-comparison}

这一维度检查谁获得服务或资源、谁承担申请和使用成本。先明确目标人群及分配单位，再比较各组的获得比例、资源数量或负担分布。例如培训机会应以符合所声明申请范围的人群为分母，而不是只统计最终完成线上报名者，否则可能遗漏在申请入口就被排除的人。

若保护要求是“各群体总体获得名额的比例应当相同”，比较时不再以资格细分人群，于是检查各组所有申请者中的获选比例。这个要求可以用\term{人口统计均等}（demographic parity）表示：
\[
 \mathbb P(\hat y=1\mid a=u)
 =\mathbb P(\hat y=1\mid a=v).
\]
它直接约束机会分配，却没有使用资格标签。若标签存在明显偏差，这种不依赖标签的特征可能有意义；若资格差异是任务需要考虑的信息，又需要解释为何仍要求比例相同。

机会比例还可能用差值或比值报告。例如两个群体录取率为$0.02$与$0.04$时，绝对差仅为两个百分点，前者却只有后者的一半。差值便于表达额外影响了多少人，比值强调相对机会，两者都不能省略基准比例。

实际检验时，应在同一任务和时间范围内记录组别、是否获得机会以及等待时间、费用等负担。录取率用组内获选人数除以组内人数，群体间可以比较差值或比值；负担则可比较平均值、分位数或超过已声明上限的比例。容许差异与上限应由任务要求确定。最终录取率相同而某组需要更多申诉时，机会与过程负担会给出不同结论。

这里对总体获选比例的比较没有直接说明差异的原因，也没有以真实资格为条件。要判断合格者是否受到不同程度的误拒，还需下一小节的标签与错误分析。对于某一大群体内部是否藏有集中伤害，则应先细化当前比较的人群。

\paragraph{交叉群体中的细化检验}

只检查几个大群体仍可能遗漏集中出现的伤害。设两个二元属性将人群均匀划成四组，四组录取率依次为$0.9,0.1,0.1,0.9$。分别按任何一个属性统计，两组录取率都为$0.5$，交叉组之间却相差$0.8$。Kearns等将这种粗粒度统计通过而结构化子群失败的现象称为公平性划分操纵，并研究对子群族同时审计和学习的方法\scite{kearns2018subgroup}

令$\mathcal G$为预先约定或受复杂度限制的子群判别函数族，
$g(\bm{x},a)=1$表示属于某一子群。要定义\term{子群公平}（subgroup fairness），
还必须固定参考总体。用事件$C$表示当前比较所针对的人群：人口统计均等可取$C$为
全部申请者，假正例率比较则取$C=\{y=0\}$，并要求$\mathbb P(C)>0$。记
\begin{align*}
 r_C&=\mathbb P(\hat y=1\mid C),\\
 \alpha_C(g)&=\mathbb P(g=1,C),\\
 r_C(g)&=\mathbb P(\hat y=1\mid g=1,C)
 \quad\text{当 }\alpha_C(g)>0.
\end{align*}
其中$r_C$是参考总体的正决策率，$r_C(g)$是子群内的对应比例，
$\alpha_C(g)$是该子群在总体中的概率质量。对子群$g$定义加权违规量
\begin{equation}
 \begin{aligned}
 v_C(g)&=
 \begin{cases}
  \alpha_C(g)\,|r_C(g)-r_C|,&\alpha_C(g)>0,\\
  0,&\alpha_C(g)=0,
 \end{cases}\\
 \sup_{g\in\mathcal G}v_C(g)&\leq\gamma,\qquad \gamma\in[0,1].
 \end{aligned}
 \label{eq:fairness-subgroup-violation}
\end{equation}
给定容差$\gamma$后，最后一个不等式才构成对整个子群族的要求。分段定义避免在零概率
子群上调用没有定义的条件概率。概率质量权重避免把极小子群中高度不稳定的原始比例差
直接当成同等规模的总体违规；它也意味着报告结果时仍应同时给出子群大小和未加权差异。

检验时，审计器在$\mathcal G$中搜索加权违规量$v_C(g)$最大的子群，再用独立数据核验
找到的差异，而不是只搜索原始率差最大的子群。这里改变的是比较人群的粒度，所比较的
仍是已经选定的机会比例或条件错误率。若参考总体、容差或子群族改变，得到的是另一个
公平要求；若子群族不包含真正受损的人群，或搜索算法找不到违规，未发现差异也不能排除
其他群体受损。

\subsection{预测错误与分数含义}
\label{sec:fairness-prediction-comparison}

获得机会的比例相同，不代表错误由各组同样承担。这一维度检查模型在不同群体上是否具有可接受的错误水平，以及同一分数是否具有相同含义。需要在预测与决策记录之外取得可信的真实标签；如果历史标签已经偏离认可的资格标准，精确计算错误率也可能检验错对象。

若重点是避免合格者因群体不同而被不同程度地拒绝，可以要求\term{机会均等}（equal opportunity），即各组真正例率相同。进一步要求真正例率和假正例率都相同，则得到\term{均等化赔率}（equalized odds）：
\begin{equation}
 \mathbb P(\hat y=1\mid y=b,a=u)
 =\mathbb P(\hat y=1\mid y=b,a=v),
 \quad b\in\{0,1\}.
 \label{eq:fairness-equalized-odds}
\end{equation}
这些标准把标签作为比较资格的依据，其适用性随标签含义而变化。Hardt等讨论了这类群体判定及给定预测器的后处理，并指出只依据联合统计量进行判断的局限\scite{hardt2016opportunity}

概率校准要求具有相同预测概率的样本呈现相应的事件频率（见第~\ref{chap:trust-reliability}章）。公平性还可以进一步要求预测分数在各组中具有相同解释。例如，对预测合格概率$p(\bm{x})$要求组内校准：
$\mathbb E[y\mid p(\bm{x}),a]=p(\bm{x})$。
它要求同一分数在各组对应相同事件频率，却不直接要求获得机会的比例或错误率相同。

这些要求可以按条件独立关系理解。人口统计均等要求$\hat y$与$a$独立，均等化赔率要求给定$y$后$\hat y$与$a$独立。若要求同一决策结果在各组对应相同的标签比例，则考察给定$\hat y$后$y$与$a$是否独立；其中对$\hat y=1$的要求，就是正预测值一致。条件放在不同位置，意味着保护对象不同：申请者关注同等资格能否获得机会，名额分配者可能关注被选者达到目标的概率。必须说明采用哪一个视角，以及为什么这一视角适合当前任务。

\paragraph{按标签与分数进行检验}

若组内没有正例，真正例率没有定义；把这种情况记为零再计算差距，会把缺少证据误写成最差表现。其他条件概率同样需要检查分母，不能只保存最后一个差距数。

\begin{example}[同一批申请者上的不同公平指标]
\label{ex:fairness-rates}
构造两个各100人的群体，其中合格者分别为20人和60人。第一组有16名合格者与8名不合格者获得机会，第二组对应为48名与4名。两组的获选比例为$24/100=0.24$与$52/100=0.52$；真正例率为$16/20=0.8$与$48/60=0.8$；假正例率为$8/80=0.1$与$4/40=0.1$；正预测值为$16/24\approx0.667$与$48/52\approx0.923$。

因此，在这批构造数据上，均等化赔率的两项差距均为零，人口统计均等的差距却为$0.28$，正预测值差距约为$0.256$。若把公平仅记为“两个群体相同”，就无法说明其中哪一个结论成立。这里给出的都是样本比例，推向未来申请者时仍需估计误差。
\end{example}

这个例子的混淆矩阵尚不能检验分数校准，因为它没有记录每个人的预测概率。假设在两组获选者中各有10人得到$0.8$分，其中实际合格人数分别为8人和6人，则这个分数组中的实际比例是$0.8$和$0.6$：第一组的经验频率与该分数一致，第二组的经验频率比该分数低$0.2$。这与上例给出的总体计数可以同时成立；是否存在总体失准，仍需结合样本量评估估计误差。分数校准检验提供了二元错误率没有包含的信息，也不能被“真正例率相同”替代。

在实际审计中，应根据选定准则保存各组的混淆矩阵与分数。机会比例用组内获选人数除以组内总人数，真正例率用合格且获选人数除以合格人数，正预测值则用合格且获选人数除以获选人数；校准检验按组和分数区间比较平均预测与实际频率。这些统计量可以来自同一批记录，却因条件和分母不同回答不同问题。将差值与事先确定的容许差异比较之前，还要检查估计误差，见第~\ref{sec:fairness-evidence}节。

\paragraph{交叉群体中的分数检验}

第~\ref{sec:fairness-group-comparison}节中细化比较人群的思路也可以用于检查概率分数。\term{多重校准}（multicalibration）要求预测在一族可识别子群的各个分数区间内，都与实际频率大致一致。Hébert-Johnson等提出这一要求，以避免总体校准掩盖复杂子群中的系统性偏差\scite{hebertjohnson2018multicalibration}

沿用第~\ref{sec:reliability-evaluation}节的分箱检验，对分数区间$B$和子群$g$可检查
\begin{equation}
 \left|\mathbb E[y-p(\bm{x})\mid
       g(\bm{x},a)=1,\ p(\bm{x})\in B]\right|
 \leq\varepsilon.
 \label{eq:fairness-multicalibration}
\end{equation}
这里还需给出参与检查的最小群体质量及区间划分，避免对几乎没有样本的条件作不可能的精确要求。实际检验把条件期望替换为对应样本中的平均残差，并为所搜索的子群和分数区间报告误差范围。这并不要求各组成功率相同，而是要求模型不要持续高估或低估它能够识别的子群。

\subsection{相似个体的处理差异}
\label{sec:fairness-individual-comparison}

群体平均可能掩盖组内差异。\term{个体公平}（individual fairness）要求任务相关条件相近的个体获得相近处理。这个要求把困难转移到“相关”和“相近”的定义：使用什么距离、哪些因素合理、哪些因素不应参与，都需要说明。

Dwork等在Fairness Through Awareness中将这一要求表达为相似性约束。设$q_{\bm{x}}$是对申请者$\bm{x}$采取各种决策的概率分布，$d(\bm{x},\bm{x}')$表示任务认可的个体差异，$D$衡量决策分布差异，则要求
\begin{equation}
 D(q_{\bm{x}},q_{\bm{x}'})\leq d(\bm{x},\bm{x}').
 \label{eq:fairness-individual}
\end{equation}
这一判定限制相近者之间的决策概率差异。二元决策若用总变差距离，左侧就是两人的录取概率之差的绝对值\scite{dwork2012awareness}

例如两人的任务距离为$0.05$，录取概率却分别为$0.9$和$0.1$，就违反了上述条件；分别为$0.60$和$0.63$则符合这一对个体的要求。这一定义比群体平均更细，但距离并非免费获得。若把历史收入差异直接写入距离，历史机会不平等可能被当成“本来就不同”。某个距离下通过检验，只能支持相对于该距离的公平结论。随机化能帮助满足平滑的相似性要求，也意味着同等个体在一次抽签中仍可能得到不同结果；这一公平约束针对的是事前分布。

检验个体公平需要先固定任务距离和决策分布的估计方式，再搜索差异过大的个体对。对每一对记录可计算$D(q_{\bm{x}},q_{\bm{x}'})-d(\bm{x},\bm{x}')$，正值就是违反式~\eqref{eq:fairness-individual}的证据。若模型直接提供随机化决策的概率，可据此计算；若只能重复调用随机化接口，则需估计这些概率及其误差。测试若只覆盖有限个体对，应报告搜索范围和最大已发现违规量，不能把未发现违规写成所有个体对均满足约束。群体平均检验与这种成对检验相互补充，前者通过并不推出后者通过。

可以用四个人看出这种不互推。两组各有一名高资格者和一名低资格者，第一组录取高资格者，第二组录取低资格者。两组录取率都为一半，但同等资格者受到相反处理，因此人口统计均等并不保证个体公平。反过来，换成高资格者比例不同的两个群体，若规则一致地录取所有高资格者，而且任务距离将同等资格者之间的距离设为0、高低资格者之间的距离设为1，式~\eqref{eq:fairness-individual}就能成立；两组高资格者比例不同，仍会得到不同录取率。个体要求依赖的是任务距离，人口统计均等依赖的是组内总体比例，两者保护的关系不同。

\subsection{决策依据与影响路径}
\label{sec:fairness-counterfactual-comparison}

若关注的是不合理因素通过何种路径影响决策，就需要进一步建立因果模型。\term{反事实公平}（counterfactual fairness）比较同一个人在不同群体属性设定下的预测，而不是拿两个现实中的人直接作比较。Kusner等的定义固定用于刻画个体的外生背景，再考察对群体属性作干预后预测分布是否改变\scite{kusner2017counterfactual}

设$\bm{u}$为因果模型中的外生变量，$\hat y_{a\leftarrow v}(\bm{u})$表示把群体属性设为$v$时得到的预测。用$a_0$表示当前个体实际观测到的群体属性值，避免与外生变量混淆。对于实际观测$(\bm{x},a=a_0)$，一种标准写法是对所有可取值$b,v$要求
\begin{equation}
\begin{aligned}
 &\mathbb P\bigl(\hat y_{a\leftarrow a_0}(\bm{u})=b
       \mid\bm{x},a=a_0\bigr)\\
 &\qquad=\mathbb P\bigl(\hat y_{a\leftarrow v}(\bm{u})=b
       \mid\bm{x},a=a_0\bigr).
\end{aligned}
 \label{eq:fairness-counterfactual}
\end{equation}
条件中的实际观测用于推断外生背景，干预后的世界再由同一因果机制生成。这个顺序很重要：只修改数据表的一列、把其后果变量全部固定，通常不等于计算该反事实。

在一个构造模型中，测验分数满足$x=u+\gamma a$，其中$u$表示需要评价的能力，$\gamma a$表示被认定不应影响分配的系统性因素。直接用$x$预测会随$a$改变；若结构式和$\gamma$正确，使用$x-\gamma a=u$就消除了这条路径。若进一步设$\gamma=2$，实际观测为$a=1,x=7$，则结构式推得$u=5$。干预为$a=0$后，同一人的分数应变为$x=5$；使用$x$作依据的规则可能改变决策，使用$x-2a$的规则在两个世界都得到5。不能只把输入中的$a$改成0而仍保留$x=7$，因为那没有沿结构式传播干预的后果。这个例子也说明，计算公平分数有时反而需要群体属性。若实际差异来自其他未观测原因，机械减去$\gamma a$便未必合理；从观测分布通常不能唯一恢复上述结构。

有些任务只禁止部分影响路径，例如允许与目标相关的实际能力影响录取，但禁止与任务无关的直接偏好。此时需要说明路径选择和识别假设，不能把“只改变群体属性后输出不变”当成适合一切任务的要求。后文第~\ref{sec:interpretability-recourse}节将讨论决策补救，它同样涉及反事实，但寻找的是个体能够实施的行动；本节比较的群体属性干预用于分析影响机制，不意味着这些属性是个体应当改变的行动。

这类判定所需的证据与群体统计、个体成对比较不同。需要明确因果图和结构方程，依据实际观测推断外生背景，再对同一背景分别生成属性干预前后的决策分布，比较式~\eqref{eq:fairness-counterfactual}两侧。因果机制无法由现有数据识别时，应报告在不同合理机制下结论是否改变；仅有一张群体混淆矩阵，不能验证反事实公平。由模型生成的反事实比较即使完全一致，结论也仍以该因果模型成立为条件。

\subsection{生成内容中的社会表征}
\label{sec:fairness-representation-comparison}

对于生成式培训建议，还需检查模型怎样描述不同人群。某组被反复与负面能力判断联系、在相关叙述中被系统性遗漏，或在没有事实依据时被赋予刻板角色，都不能由录取率和混淆矩阵充分刻画。应先明确要评价的表达及其语境，例如“证据不足时不按群体身份推断能力”，再建立可复核的内容标注规则。

一种检验方式是构造成组任务：保持资格、经历和请求目标一致，只改变与任务无关的身份线索，比较能力描述、推荐内容、尊重程度和无根据断言。同时保留证据充分与不充分两类条件，避免模型通过一律拒答获得表面一致。需要核查成组文本是否同样自然、是否改变了事实含义，再由明确的评分规则和必要的人工复核评价输出。更换身份词的文本比较只能支持相应行为差异，不能直接等同于前一小节中的个体因果反事实。

生成具有随机性时，应重复采样并报告各类输出的比例及其波动，检查不同提示与表达下结论是否保持。内容评价者或自动评分器本身也可能有偏差，需要用经过核查的样本检查评分一致性及典型误判。第~\ref{sec:fairness-frontiers}节以BBQ等工作进一步讨论如何把证据条件落实为生成任务的评价设计。

\subsection{公平要求的声明}
\label{sec:fairness-requirements}

形成判定结论之前，应把“希望更公平”写成能够检查的要求。声明需要交代适用任务与人群、保护哪些群体或个体关系、认可什么资格标准、禁止哪些影响路径，以及优先改善哪一种后果。随后指定主要判定准则、容许差异、总体效用和资源限制，并说明采用哪些数据、在哪个时间范围内验证。这些内容共同决定怎样解释前述证据，单独给出一个公平指标还不是完整要求。

例如，可以为教学用培训分配系统声明：以独立核验的资格为$y$，对事先列出的两个群体优先限制合格者获选率的绝对差不超过0.05，同时要求两组获选率在合格者中都不低于0.8；名额总量不能超过预算，并同步报告假正例率、交叉群体差异及申请成本。数值只是教学设定，不是通用标准。差距约束限制不一致程度，最低水平要求约束绝对表现，在某组尚低于下限时，不能仅靠降低表现较好组的机会达标；若资源与这些要求不能同时满足，需要重新讨论目标或资源，不能隐去其中一项。

声明还应规定信息如何使用。例如，群体属性可在训练和审计时用于构造约束，但部署时是否允许读取该属性，要由任务要求明确决定。这会直接影响可选方法：组别阈值需要在决策时识别组别，而训练时采用公平约束并不必然要求部署模型读取群体属性。

目标声明与达成声明必须分开。前者说明系统打算满足什么，后者报告哪个模型版本在什么人群、数据和假设下获得了何种证据。独立检验未完成、关键群体样本不足或因果机制无法识别时，应保留这些限制，不能因为采用了某种去偏算法就宣称“模型公平”。

\subsection{判定结论的形成}
\label{sec:fairness-compatibility}
\label{sec:fairness-evidence}

前五个方面提供了针对不同要求的证据。形成结论时，需要明确本次采用哪些要求、它们能否同时成立，以及现有样本和假设支持多强的判断。某个维度达标只能支持限定范围的结论，不应汇总为没有条件的“系统公平”。

\paragraph{准则之间的取舍}

同一决策系统可能同时接受多种要求。出现指标分歧时，应先判断分歧是否来自准则本身，而不立即归咎于训练不足。下面在群体比较中给出几项可直接推导的冲突，它们约束的是能够共同选择的目标，不是检验公平性的额外方法。

例~\ref{ex:fairness-rates}已经显示同一批预测可以通过一种准则而不通过另一种。要区分这是偶然数值还是一般限制，可以从各组正例比例推导。记组内正例比例为$\pi_a$、真正例率为$t_a$、假正例率为$f_a$，由对真实资格分类求和，获选比例为
\[
 r_a=\pi_a t_a+(1-\pi_a)f_a.
\]
即使固定$t_a=t$与$f_a=f$，只要$t\ne f$且组间$\pi_a$不同，机会比例仍然不同。若又要求$r_a$相同，就不能任意保留前述两个共同错误率。

对获选者反过来询问资格比例，会得到另一项限制。设各组满足相同的真正例率$t$和假正例率$f$，组$a$中的正例比例为$\pi_a\in(0,1)$，则获得正决策者中的真实正例比例为
\begin{equation}
 \mathbb P(y=1\mid\hat y=1,a)
 =\frac{t\pi_a}{t\pi_a+f(1-\pi_a)}.
 \label{eq:fairness-predictive-value}
\end{equation}
当$t>0$且$f>0$时，右侧随$\pi_a$严格增加。因而基准比例不同的组，即使错误率已经相同，正决策的预测值仍然不同。这一具体结论不要求模型采用某种算法，也不适用于所有退化边界；例如$f=0$时正预测值可能均为1。关于风险分数不同公平条件之间的更一般取舍，需保留相应条件理解\scite{kleinberg2016fairness}

分数校准与错误平衡的冲突也能直接推导。设$S=p(\bm{x})\in[0,1]$在每组校准，$\mu_{1,a}=\mathbb E[S\mid y=1,a]$、$\mu_{0,a}=\mathbb E[S\mid y=0,a]$。若额外要求两类人的平均分数分别跨组相同，记为$\mu_1$和$\mu_0$，则全期望给出
\[
 \pi_a=\mathbb E[S\mid a]
       =\pi_a\mu_1+(1-\pi_a)\mu_0.
\]
两个组的$\pi_a$不同时，相减得到$\mu_1-\mu_0=1$。由于分数在$[0,1]$内，只能有$\mu_1=1,\mu_0=0$，即正例几乎必然得1分、负例几乎必然得0分。这说明在不同基准比例、非完美预测的情形下，三项要求不能同时成立。这里的“平均分数平衡”不同于式~\eqref{eq:fairness-equalized-odds}中的二元错误率相同，不能把两个论证的条件混在一起。

不可兼得关系并不意味着公平性无从讨论，而是要求在冲突出现之前说明优先保护什么。例如培训机会应优先避免遗漏合格申请者时，可以把机会均等作为主要约束，同时公开假正例率、名额总量和预测值的变化。此时没有选择的准则仍应报告，用来揭示代价；不能因为它不再是优化目标，就把相应差异从评价中删去。

\paragraph{有限样本与建模假设}

公平准则给出了要检查的关系，实际数据只能提供有限证据。群体检验依赖样本覆盖和分母，个体检验还依赖所接受的距离，反事实检验则额外依赖因果机制。后两类假设即使样本增加也不会自动正确；应把假设敏感性与抽样误差分别报告。以下以群体差距为例说明怎样从估计值形成有边界的判定。

估计群体指标时应报告分母与样本量。机会均等的分母是组内正例数，假正例率的分母是组内负例数。小群体或交叉群体可能缺少足够样本；同时搜索许多群体又会增加偶然差异。敏感属性缺失时，补全属性的误差也会进入公平评价，不能将推断群体当作无误标签。

样本误差与实质差异是不同问题。设两个固定群体各有$n_u,n_v$个独立正例，其真正例率估计为$\hat t_u,\hat t_v$。将Hoeffding界分别用于两组再取联合界，以至少$1-\delta$的概率有
\[
 |(t_u-t_v)-(\hat t_u-\hat t_v)|
 \leq
 \sqrt{\frac{\log(4/\delta)}{2n_u}}
 +\sqrt{\frac{\log(4/\delta)}{2n_v}}.
\]
即使总体样本很多，只要其中一组正例太少，差距的不确定性仍会很大。若要求绝对差距$|t_u-t_v|$不超过容许值$\eta$，应将经验绝对差$|\hat t_u-\hat t_v|$加上上式右侧的误差界，再检查所得上界是否不超过$\eta$；“没有拒绝零差异”并不能证明实际差异很小。若群体是在同一数据上反复搜索出来的，还需单独验证或控制整个搜索族的误差。

\paragraph{达标、违反与证据不足}

对于已声明的差异上限，应根据估计范围区分三种结论：差异的可靠上界仍不超过容许值，支持在相应条件下达标；可靠下界已经超过容许值，支持存在违反；估计范围跨过界限，则证据不足。个体比较还要报告距离定义和搜索范围，因果比较要报告机制假设，内容比较要报告任务与评分条件。不能把“未找到违规”或“差异不显著”直接写成已经公平。

判定报告应保留各组原始水平、使用的准则、检验数据与模型版本，并说明未覆盖的人群和未核实的假设。这里得到的是干预的目标与证据：哪些要求没有满足、在哪些人群或条件下失败。至于应补充数据、重训模型还是调整服务流程，还需结合前节的形成原因决定。

\section{公平性的干预}
\label{sec:fairness-intervention}

公平判定指出哪些要求没有满足，干预要进一步确定应该改变哪个环节。同一机会差距可能来自群体数据缺失、标签不合理、训练时忽视少数样本，也可能来自分数到行动的规则。因此，应依据已经声明的要求和形成原因，选择数据、训练或决策流程上的措施，不能看到差距便直接更换算法。

前一节已经确定保护对象、准则、容许差异和证据边界。下面三类实施方向分别回答：学习依据是否需要修正，模型学到的规则是否需要调整，以及既有分数怎样转化为恰当服务。它们可以组合。例如补齐小群体数据后仍需调整错误成本；训练约束达标后仍需检查等待与申诉负担。效果验证与持续调整则检查这些措施是否真正改善了目标人群的处境。

\subsection{数据采集与质量的改进}
\label{sec:fairness-data-intervention}

若诊断表明问题来自看不见某些人，干预重点是扩展采集渠道与覆盖范围。可以按目标人群、业务条件和交叉群体检查缺口，对缺少的语言、设备、地区或资格区间有目的地补充样本。采集“更全面”应落实为哪些条件此前缺少证据、现在增加了多少有效观测，而不只是把每组人数做成相同。重复抽取旧样本会改变训练频率，却没有带来新的个体信息。

补充数据还需区分训练与评价职责。训练可以为困难或少数群体增加样本，评价则要说明结果对应目标人群还是刻意平衡后的审计人群。若采样比例改变，直接平均会改变总体指标的含义；应保留分组结果，必要时按已知目标构成汇总，并说明权重依据。新的审计样本应尽量独立于用于选择干预方法的数据，避免把反复调试后的效果当作独立证据。

若问题来自测量，应检查同一字段在不同群体中是否表达同一种任务信息。可以对不同语言或设备条件进行配对核验，与更可靠的测量参照比较，区分真实差异与测量误差。若问题来自标签，则需要重新核对资格标准、标注指南和争议样本，必要时重新标注。不能为使群体比例一致而任意改写真实结果；要改变的是错误记录或不合理的目标定义，并保留变更依据。

标签缺失又是另一种情形。对只有获选者才有后续成绩的任务，不能把未获选者记为失败；可以寻找独立资格测验、适当的后续调查或其他有依据的观测来源。若采用改变选择机制的方式获取新证据，需要说明实施条件；依靠统计校正则要核查选择概率与识别假设，具体限制见第~\ref{sec:fairness-long-term}节。增加已有获选者的记录并不能自动填补被拒者的信息缺口。

数据修正应留下可追溯记录。Gebru等提出的数据集说明书记录数据的用途、构成、采集过程和适用范围，为检查遗漏人群及不适当使用提供依据。文档帮助定位问题，但不能替代实际补采、测量核验或标签修正\scite{gebru2021datasheets}

数据质量改善为公平学习提供更好的依据，却不自动决定资源应如何分配。若标签和覆盖已经合理，而训练仍允许牺牲小群体来换取总体收益，就需要改变学习规则；若问题发生在分数转为行动的环节，则要继续检查决策流程。

\subsection{训练目标与学习规则的调整}
\label{sec:fairness-learning-intervention}

当问题来自模型怎样使用已有数据时，可以改变不同错误在训练中的代价、限制表示中传递的信息，或直接约束所选公平指标。这些措施共同改变学得的预测规则，但作用对象不同。重加权照顾哪些样本被目标函数重视，表示约束限制怎样编码依据，显式公平约束则直接反馈某项要求的违规量。

学习阶段可以将公平要求写为约束。设$R(f)$为任务风险，$\Delta(f)$为已经选定的公平差异指标，$\eta$为容许差异，则求解
\[
 \min_f R(f),\qquad \Delta(f)\leq\eta
\]
明确了任务目标与限制。使用$R(f)+\lambda\Delta(f)$的惩罚形式便于优化，但给定$\lambda$不自动保证约束达标；仍需在独立数据上验证。

\paragraph{不同样本与错误的训练代价}

重加权改变的是训练目标中每组的相对贡献。假设训练集中两组样本数为$n_u,n_v$，总数$n=n_u+n_v$，希望每组对经验损失各占一半，可以将组$a$的每个样本赋权$\omega_a=n/(2n_a)$。此时
\[
 \frac1n\sum_{i=1}^{n}\omega_{a_i}\ell(f(\bm x_i),y_i)
 =\frac12\hat R_u(f)+\frac12\hat R_v(f),
\]
其中$\hat R_a$是组内平均损失。实际做法是根据训练集确定权重，在每次梯度计算中按权重累加样本损失，再在独立数据上重新计算已选公平指标。它让小群体不会只因人数少而被目标函数忽视，却不自动使两组错误率相同，也不修正错误标签。是否采用各占一半的目标仍须由任务要求决定。

\paragraph{表示中的信息依赖}

若同一套特征将被多个预测器使用，可以在表示层减少不希望传递的群体信息。Zemel等的Learning Fair Representations用软分配把样本映射到一组原型，同时权衡输入重构、标签预测和不同群体的原型分配差异。这里不是简单删除一个字段，而是检查变换后的表示是否仍以相近方式编码不同群体\scite{zemel2013fairrepresentations}

设$q_k(\bm{x})$为输入分配给第$k$个原型的概率，$\bar q_{k,a}$为组内平均分配概率，一个便于理解的公平惩罚是$\sum_k|\bar q_{k,u}-\bar q_{k,v}|$。如果之后先按$q_k$抽取一个原型，再仅根据原型作决策，且每个原型的组内平均分配相同，那么由全概率公式可以得到相同录取率。但有限样本上的近似相同，或者直接将整个软分配向量交给任意下游模型，都不能无条件继承这个结论。表示经过不同使用方式后，需要重新检查保护目标。

这类表示干预在训练时需要同时调整原型、样本到原型的分配以及基于表示的任务预测器。输入重构损失限制“丢掉所有信息”的退化解，标签预测损失保留任务所需区分能力，群体分配惩罚则减少指定的信息差异。三项权重确定取舍，训练后在独立数据上分别检查任务效用、表示中的群体可预测性和最终决策的公平指标；仅看到训练时分配惩罚减小，还不能断言下游使用已经公平。

另一条路线是让一个辅助预测器主动寻找残留的群体信息。\term{对抗去偏}（adversarial debiasing）在训练任务预测器时，同时训练一个根据其输出或中间表示预测群体属性的对手。Zhang等将这种相互竞争用于降低特定公平差异；其作用取决于对手看见什么信息，以及它是否足够强\scite{zhang2018debiasing}

若用$h_{\bm{\theta}}$表示待学习输出，$d_{\bm{\phi}}$表示群体预测器，一个概括这种竞争的目标是
\[
 \min_{\bm{\theta}}\max_{\bm{\phi}}
 \left\{L_{\mathrm{task}}(\bm{\theta})
       -\lambda L_{\mathrm{group}}(\bm{\theta},\bm{\phi})\right\}.
\]
对手通过减小群体预测损失寻找泄漏，任务模型则试图增大对手的最优损失。训练时交替进行两类更新：固定任务模型，用真实群体属性训练对手；再固定对手，以降低任务损失、增大群体预测损失的方向更新任务模型。部署后通常只保留任务模型。评价时还应另训独立的群体预测器，检查信息是否只是绕过了训练中的对手。只让对手观察输出，对应消除输出与群体的依赖；若追求均等化赔率，需要处理给定真实标签后的剩余依赖，例如在各标签层内训练对手。对手未超过基准预测水平，只能支持相对于其能力与训练过程的判断，不能证明不存在更强的群体推断器。对抗优化不稳定、任务信息与群体信息重叠，都可能使准确性与公平目标发生取舍。

\paragraph{公平约束下的分类学习}

表示学习通过改变信息间接影响公平性；若保护标准已经明确，还可以直接优化相应条件矩。Agarwal等把公平分类归约为一系列\term{成本敏感分类}（cost-sensitive classification）问题，即为不同样本的不同错误赋予不同代价，再调用已有分类学习器。归约允许学习器保持自身表示形式，而把公平要求放在外层求解\scite{agarwal2018fairreductions}

将选定的经验公平约束拆成条件矩差的单向限制，写为$\hat g_j(f)\leq\eta_j$，这样约束量对分类器的随机混合保持线性。对分类器上的随机混合分布$Q$，可构造
\begin{equation}
\begin{aligned}
 &\min_Q\sup_{\bm{\lambda}\geq\bm{0}}
 \left\{\mathbb E_{f\sim Q}[\hat R(f)]\right.\\
 &\qquad\left.+
 \sum_j\lambda_j
 \bigl(\mathbb E_{f\sim Q}[\hat g_j(f)]-\eta_j\bigr)
 \right\}.
\end{aligned}
 \label{eq:fairness-reduction}
\end{equation}
这是解释约束作用的拉格朗日形式；具体算法还会限定乘子的搜索范围与停止精度。乘子提高某项违规的代价，分类学习器则在调整后的成本下寻找更好的预测器。多轮得到的分类器可以混合，使无法由单个确定规则达到的中间错误率成为可达目标。

以提高某组真正例率为例，外层可以增大该组正例被误拒的成本，使学习器更愿意将这些样本判为正例。若这又明显增加假正例率，另一项约束的乘子会抑制这一变化。这个过程比“公平惩罚越大越好”更具体：每项约束对应哪些样本、哪些错误，可以追踪。

实施时，可从一组初始乘子出发训练成本敏感分类器，计算每项经验约束的违规量，再增大被违反约束的乘子，重复学习并保存所得分类器。最终按求得的混合权重组成规则，而不是只挑训练准确率最高的一轮。这里直接反馈的是保护准则的违规量；前述对抗去偏反馈的是群体信息是否仍可预测。两者有时都能减少某项群体差距，却不能仅凭采用了同一种优化形式便视为相同方法。

归约的优化性质以成本敏感学习器的求解能力、候选分类器族和允许的优化误差为条件。经验约束达标也不等于总体约束必然达标，仍需独立评价和适当容差。群体属性可以只用于训练时构造约束，部署时分类器使用哪些特征由所选模型族决定；这与必须在使用时读取群体属性的组别后处理有所不同。

当保护对象包括大量交叉群体时，无法总在训练前列完每一项约束。Kearns等的学习与审计交替方法让审计器搜索当前分类器加权违规量最大的子群，再由学习器针对反馈修改分类器。这与第~\ref{sec:fairness-group-comparison}节的审计使用同一种参考总体、概率质量和差异证据，但这里把发现的违规转为下一轮学习约束。每轮改善一个子群可能影响另一个交叠子群，停止时仍须检查约定的整个群体族，并在未参与迭代的数据上验证。

若声明的是个体公平，则可以在任务距离明确的条件下，对相近个体的决策分布增加一致性约束，或惩罚超过允许距离的差异；训练后仍要用未参与优化的个体对检验。若声明的是反事实公平，则需要依据已认可的因果模型构造不受被禁止路径影响的预测依据，或约束相应反事实世界中的预测一致性。前述测验例子中使用$x-\gamma a$就是在特定结构成立时构造依据的一种方式；机制假设不可靠时，单纯优化一致性不能使因果结论可靠。这两类要求不应被总体群体比例相等替代。

\subsection{决策规则与服务流程的调整}
\label{sec:fairness-decision-intervention}

模型无法重训时，仍需区分修改的对象：分数失准可以在输出分数上作修正，机会与错误率不符合约束则需要改变从分数到行动的规则。这两种操作可能相互影响，但一个目标达标并不推出另一个目标达标。

针对式~\eqref{eq:fairness-multicalibration}中的失准，一个直观的多重校准修正过程是在“子群与分数区间”中按平均残差调整预测，再检查其他交叠群体。正式算法还要控制调整步长、搜索代价和泛化误差。分数修正后必须重新检验各组校准；若下游使用固定阈值，修正还可能改变谁获得机会，因此也要重算相应的分配指标。

如果保留原分数，只修改决策规则，可以按群体选择阈值。对每个群体改变阈值会得到一系列假正例率与真正例率坐标，即ROC曲线上的点。随机选择两个阈值，对应两点的凸组合。Hardt等的后处理利用这些可达区域，在满足均等化赔率的共同坐标中选择合适的误差水平。

\begin{example}[随机化如何补出共同错误率]
\label{ex:fairness-random-threshold}
假设在独立验证数据上，第一组的两个阈值分别得到$(f,t)=(0.1,0.5)$和$(0.3,0.9)$，第二组某个阈值得到$(0.2,0.7)$。对第一组以各一半概率使用两个阈值，期望错误率成为$(0.2,0.7)$，从而与第二组一致。随机化必须独立于尚未观察的真实标签；部署后仍需检查这些验证率能否代表实际分布。
\end{example}

如果两组共同可达的区域很差，后处理可能主要通过降低较好组表现实现均衡。它不能创造原分数中没有的区分信息，也不保证满足个体公平：两个任务条件相近的人，可能因为组别阈值不同而获得不同录取概率。后处理的优点是可清楚计算固定分数上的取舍，限制则是同时固定了已有表示与排序。是否使用这种规则，必须结合允许的信息使用、资源约束和受影响者的实际处境判断。

决策干预还包括模型之外的服务过程。对模型难以可靠处理的申请，可以设置复核、补充材料与纠错入口，并记录处理标准和完成时间；对申请渠道本身的障碍，可以改善语言支持或材料提交方式。人工参与并不自动保证公平，仍要检查不同群体进入复核的比例、等待负担和最终结果。若某组被大量转交人工而得不到及时服务，就不能只因自动预测中的错误减少而宣称改善。

对于生成式建议，可以在输出与服务规则中明确应依据什么事实、何时请求补充信息、哪些表述需要复核，并在前述成组任务上检查这些调整的效果。训练阶段也可利用经核查的对照样本或反馈调整行为，但“使用了去偏数据”与“已满足内容公平要求”仍是两个不同结论。

\subsection{干预效果与持续调整}
\label{sec:fairness-intervention-evaluation}
\label{sec:fairness-long-term}

数据、训练与流程上的措施都必须回到同一份公平要求接受检验。即时检验比较修改前后是否改善，长期跟踪则检查干预改变数据与人群状态后是否仍然有效；二者共同决定是否保留、调整或撤回措施。

\paragraph{即时效果与负担转移}

干预效果需要相对于明确的对照判断。应固定任务、保护群体、标签定义和主要公平准则，在未用于选择措施的数据上，用同一批样本评价干预前后的模型或规则；若干预已经改变人群和标签生成过程，就需要后文讨论的额外识别条件。报告应同时包括总体效用、各组指标和保护目标，并给出变化量的误差范围。例如，机会均等的差距从$0.2$降为0，可能来自较差组的真正例率由$0.6$提高到$0.8$，也可能来自较好组由$0.8$降为$0.6$。仅报告差距会把这两个结果写成同样的改善。应同时保存两组原值、新值、假正例率和名额数量，再判断变化是否符合目标。还应检查是否把错误负担转移到另一个群体或另一个环节。

对拒答系统，需要把第~\ref{sec:reliability-evaluation}节中的自动处理比例按群体报告，并与各组接受部分的风险一起比较。若某组的自动预测错误率降低，只因为该组大量请求被拒绝，结论就不能写成“服务更加公平”。拒答后的等待、人工成本和结果也属于干预效果；这里的服务覆盖是获得处理的比例，不是预测集合包含真实标签的覆盖率。

选择措施与报告效果应使用职责明确的数据。若反复选择最有利的公平指标，容易只展示局部改善。主要标准应由使用要求确定，其他指标用于揭示代价和边界。效果不确定时，应明确样本限制，而不是把差异不显著解释为已经消除不公平。

\paragraph{选择性标签与后续结果}

即时评价固定比较条件，检查干预当时改善了什么。长期评价还要考虑干预改变条件本身：一次分配会改变下一轮可见的数据和资源。推荐系统增加某些内容的曝光后，会获得更多点击证据；培训名额分配提高某些人的能力后，又改变未来的资格分布。因此，即使静态指标改善，也需要检查长期机会、能力与负担的演化。

在选择性标签条件下，可以记录选择概率，并在条件可交换、正概率覆盖等假设下使用相应加权估计；若从未给某类申请者机会，就无法仅靠这种加权恢复其结果。需要有目的地补充证据，并检查数据获取方式是否符合使用要求。

Lakkaraju等系统研究了选择性标签给预测评价带来的困难，并利用不同决策者的选择差异构造评价证据。其意义在于指出，验证集的标签生成过程本身可能与待评价决策有关；决策者差异也只有在相应比较假设下才能支持推断，不能被自动视为随机实验\scite{lakkaraju2017selective}

为区分即时资格与后续收益，本段用$z$表示培训后的结果，而不再沿用前文的资格标签$y$。令$r=1$表示获得培训，$z(1)$表示接受培训时的潜在结果，$e(\bm{x},a)=\mathbb P(r=1\mid\bm{x},a)$。若给定$(\bm{x},a)$后，选择$r$与$z(1)$独立，且目标人群中$e>0$，还满足接受培训时观测结果$z=z(1)$，则
\begin{equation}
 \mathbb E\!\left[\frac{r z}{e(\bm{x},a)}\mid a\right]
 =\mathbb E[z(1)\mid a].
 \label{eq:fairness-selective-weight}
\end{equation}
左侧的乘积在$r=0$时定义为零，不需要观测被拒者的$z$。证明只需先条件于$(\bm{x},a)$：选择独立性使加权项的期望等于该条件下的潜在结果期望，再取组内平均。若审核者使用了数据中没有记录的重要信息，独立性就可能失败；若选择概率很小，权重会放大估计方差。这些限制比“做一次重加权”更决定结论是否可信。

长期问题还不止于标签缺失。Liu等的Delayed Impact of Fair Machine Learning用一个决策后状态变化模型分析群体受到的后续影响，表明在其模型假设下，静态公平约束可能帮助也可能损害原本希望保护的群体。决定影响方向的是谁被新增选中、选择如何改变其未来状态，以及群体原有分布\scite{liu2018delayedimpact}

可以用一个明确构造的例子理解这种机制。若对某类申请者提供一项机会，成功后状态增加1，失败后因成本增加而减少2，成功概率为$p$，则接受机会的期望变化是$p-2(1-p)=3p-2$。当$p=0.5$时，扩大该类人的机会数量却使每名新增参与者的期望状态变化为$-0.5$。这不是现实培训的经验结论，而是在说明：只有新增正决策在目标人群和当前决策规则下具有正的条件期望净收益时，“多给机会”才会提高这个模型中的群体平均状态。若要进一步声称每个人都会改善，还需要更强的个体层条件。提供费用补助、调整课程或降低失败代价，可以改变这一结果，不必只在录取阈值上作选择。

持续干预应由已经说明的标准和新证据驱动。出现新的群体差异时，应回到形成机制判断是数据变化、测量变化还是干预副作用，再调整措施，而不是只对当前差距继续加大惩罚。

\section{前沿研究}
\label{sec:fairness-frontiers}

前文已建立分类、生成内容及服务过程的基本判定方法。面向生成模型与连续服务，相关研究进一步考察三个问题：开放文本中的伤害如何转化为可检验证据，不同指标的结论是否稳定，以及一次输出之外的完整服务是否公平。它们分别扩展判定对象、测量可信度和影响范围。

\subsection{生成任务的证据条件与刻板印象}

生成内容的成组检验需要保持资格、经历与请求目标一致，只改变无关身份线索，基本方法见第~\ref{sec:fairness-representation-comparison}节。进一步构造评价时，还要控制模型究竟掌握了多少事实：缺少依据时是否诉诸刻板印象，与已有依据时是否忽略事实，是两种不同的失败。

Parrish等提出的BBQ问答基准把公平评价与证据条件联系起来：信息不足时，模型是否用群体刻板印象代替未知；信息充分时，模型是否仍让刻板印象压过明确事实。基准使用成组构造的问题控制这些条件，覆盖美国英语情境中的九类社会维度。它提供了可比较的评价设计，但其文化和语言范围也限制了结论可以推广到哪里\scite{parrish2022bbq}

例如在一个自构教学问题中，只知道两名来自不同地区的申请者都参加了测验，却不知道成绩，询问谁更适合培训时，模型应承认信息不足。若补充其中一人的成绩已经满足明确资格标准，模型又不应只因地区身份而忽略这条证据。两个条件共同区分“减少无根据的偏见”与“拒绝使用已经充分的事实”。仅靠总是回答“不知道”，可能降低某类偏见计数，却没有获得有用且公平的问答能力。

成组比较也不是自动有效的因果实验。更换身份词可能改变代词指代、语法、职业背景或文本自然程度，导致比较同时改变多个因素。评价需要人工核查语义等价性，报告事实性任务表现与群体差异，并区分模型主动生成的内容、提供候选答案时的选择以及回答拒绝。不同接口观察的是不同的行为，不能把一种接口上的改善直接写成整个模型已经消除偏见。

\subsection{评价结论对指标选择的敏感性}

即使保持同一模型，不同指标也可能给出不同排序。Asgari等在八个大语言模型、七种偏见指标和九个语料上的研究，用指标一致性分数与模型一致性分数分析这种差异，观察到一些同类指标也未表现出预期的一致关系。其结果支持在这组模型和语料中检查评价结论对指标选择的敏感性，而不是认为存在一个已经确定的通用偏见排名\scite{asgari2026biasagreement}

指标不一致至少可能来自两种原因：一种是估计噪声、提示变化或测试数据差异，另一种是指标确实关注不同伤害。毒性表达减少，不能推出职业刻板印象减少；文本概率偏好减弱，也不必然改变最终生成。前一种情况需要提高测量稳定性，后一种情况需要分别报告并说明保护目标。直接把若干量纲不同的分数求平均，会隐藏分歧来源，也可能掩盖某个群体上的严重失败。

\subsection{连续服务中的负担与结果}

输出文本的公平评价仍没有覆盖用户完成整项任务的经历。长期交互还使公平性与适应策略相连。个性化可以提高服务效果，也可能使某些群体更难获得纠正模型的机会。未来评价需要把输出内容、实际服务覆盖和反馈后的影响连接起来，并区分技术上可测量的差异与需要共同决定的价值标准。

对于代办申请或提供培训建议的智能体，一次回答只是服务链中的一步。即使推荐文本相近，不同用户可能面对不同的澄清轮次、等待时间和转人工概率。可以沿用本章的判定思路，先定义完整任务的成功与负担，再按群体分析每一步的流失及最终结果。这里与第~\ref{chap:trust-alignment}章的行为监督相接：监督不只检查是否出现禁止性内容，还要检查系统是否持续把更多服务成本转移给某些用户。

\section{本章小结}
\label{sec:fairness-summary}

回到培训名额分配，公平性需要先说明什么差异有合理依据，再检查机会、错误、个体待遇、影响路径与社会表征是否符合要求。群体比例、个体成对比较、因果反事实与生成内容评价分别提供不同证据；指标之间可能冲突，有限样本及距离、因果和评分假设又限制结论强度。因此，应给出限定任务与人群的达标、违反或证据不足结论，而不能用一个总分替代全部判断。

公平要求确定后，还需要追查不合理差异发生在哪个环节。遗漏人群需要补充覆盖，测量与标签偏差需要修正依据，训练目标的偏重需要调整学习规则，阈值与服务负担则需要改变决策流程。这些措施能够组合，却不能互相替代。最终要同时检查各组原有水平、干预后的实际收益与长期反馈，避免通过降低服务或转移负担获得表面均等。公平干预也揭示了一个更一般的问题：即使某项指标已经改善，系统仍可能没有实现它原本代表的要求。下一章将从公平约束扩展到更一般的任务目标和行为边界，讨论怎样监督、保持和干预模型行为。采集与审计所需的群体信息则还要受到隐私保护，其边界见第~\ref{chap:trust-privacy}章。

\paragraph{延伸阅读}

Hardt等关于机会均等与后处理的工作、Dwork等关于个体公平的工作，分别从群体条件率和任务相似性刻画保护要求。两者适合与本章的指标冲突一起阅读，以辨认不同准则保护的对象及其不可互推之处\scite{hardt2016opportunity,dwork2012awareness}

若关注决策依据和长期后果，可继续阅读Kusner等人的反事实公平与Liu等人的延迟影响分析。前者依赖因果模型刻画允许路径，后者说明静态约束的效果取决于决策怎样改变后续状态；两者都要求把建模假设与经验统计分开\scite{kusner2017counterfactual,liu2018delayedimpact}
